%\documentclass{beamer}
\documentclass[xcolor=svgnames,10pt]{beamer}
%%%%%%%%%%%%%%%%%%%%%%%%%%%%%%%%%%%%%%%%%%%%%%%%%%%%%%%
%%%%%%%%%%%%%%%%%%%%%%%% packages %%%%%%%%%%%%%%%%%%%%%
%\usepackage{amsmath}
%%%%%%%%%%%%%%%%%%%%%%%%%%%%%%%%%%%%%%%%%%%%%%%%%%%%%%%
%%%%%%%%%%%%%%%%%%%%%%%% theme %%%%%%%%%%%%%%%%%%%%%%%%
\usetheme{AnnArbor}  %%  in this section for example: (default, Bergen, Boadilla, Madrid, AnnArbor, CambridgUS, Pittsburgh,, Rochester, Without navigation bar If you want to have a navigation bar (Antibes, JuanLesPins, Montpellier, for branch form), (Berkeley, PaloAlto, Goettingen, Marburg, Hannover, for table of content side bar), (Berlin, Ilmenau, Dresden, Darmstadt, Frankfurt, Singapore, Szeged, for mini frame navigation), (  Copenhagen, Luebeck, Malmoe, Warsaw,  for section and subsection table)
%
\usecolortheme{crane}  %% albatross, beaver, beetle, crane, default, dolphin, dove, fly, lily, orchid, rose, seagull, seahorse, sidebartab, structure, whale, wolverine
%
\setbeamercolor{normal text}{fg=MidnightBlue}  %% font color
\setbeamercolor{normal text}{bg=yellow!20} %% background color
 %%%%%%%%%%%%%%%%% background picture%%%%%%%%%%%%%%%%%%%%%%%%%%%%%%
%\usebackgroundtemplate 
%{%
%	\includegraphics[width=\paperwidth,height=\paperheight]{back5.jpg}%
%}
%%%%%%%%%%%%%%%%%%%%%%%%%%%%%%%%%%%%%%%%%%%%%%%%%%%%%%%%%%%%%%%%%
%%%%%%%%%%%%%%%%%%%%%%%% xepersian %%%%%%%%%%%%%%%%%%%%%%%%%%%%%%
\usepackage{xepersian}
\settextfont{Yas.ttf}
%\setmathdigitfont{B Mitra}
\defpersianfont\nas[Scale=1.7]{IranNastaliq.ttf}
\defpersianfont\titr{BTitr.ttf}
%\defpersianfont\koodak{B Koodak}
%\defpersianfont\homa[Scale=1.2]{B Homa}
%\providetranslation{lemma}{لم}
%%%%%%%%%%%%%%%%%%%% Path of files %%%%%%%%%%%%%%%%%%%%%
\graphicspath{{./Fig/}}
%%%%%%%%%%%%%%%%%%%%%  Title %%%%%%%%%%%%%%%%%%%%%%%%%%%
\title{عنوان اسلاید}
\subtitle{زیرعنوان اسلاید}
\author{نام نویسنده}
\institute{نام موسسه}
\date{تاریخ}
%%%%%%%%%%%%%%%%%%%%%% Begin Document %%%%%%%%%%%%%%%%%%
\begin{document}
%%%%%%%%%%%%%%%%%%%%%% Begin Frames %%%%%%%%%%%%%%%%%%%%
{\usebackgroundtemplate{\includegraphics[width=\paperwidth,height=\paperheight]{back5.jpg}}
\begin{frame}[plain,noframenumbering]

%برای عدم نمایش نوار بالایی و پایین در فرم آپشن plain و برای اینکه جزء شماره اسلاید محسوب نشود آپشن فرم noframenumbering قرار داده شده است.
\centerline{\includegraphics<1>[width=0.5\paperwidth]{besm2-trans}}
\end{frame}
}
%%%%%%%%%%%%%%%%%%%%%%%%%%%%%%%%%%%%%%%%%%%%%%%%%%%%%%%%%%%%%
\begin{frame}
\transboxout[duration=2]
%\maketitle
\titlepage
\end{frame}
%%%%%%%%%%%%%%%%%%%%%%%%%%%%%%%%%%%%%%%%%%%%%%%%%%%%%%%
%%%%%%%%%%%%%%%%%%%%%%%%%%%%%%%%%%%%%%%%%%%%%%%%%%%%%%%%%%%%%%%%%%%%%
%\usebackgroundtemplate{\includegraphics[width=\paperwidth,height=\paperheight]{Back5}%{Pj}%{Teal}%{Wave}
 %}
 
\begin{frame}
\begin{picture}(0,0)
\put (0,-30){\centerline{\includegraphics[width=1.2cm]{YU-Logo-Ipe-6-trans}}}
\end{picture}
\vskip 0.1 in
\maketitle
\end{frame}
\section{}
%%%%%%%%%%%%%%%%%%%%%%%%%%%%%%%%%%%%%%%%%%%%%%%%%%%%%%%%%%%%%%%%%%%%%%
\section{بخش اول}
\subsection{محیط ریاضی}
%%%%%%%%%%%%%%%%%%%%%%%%%%%%%%%%%%%%%%%%%%%%%%%%%%%%%%%%%%%%%%%%%%%%%%
\begin{frame}[t]{{\titr عنوان اسلاید با فونت تیتر}}
    \begin{flushright}
        در این اسلاید محیط های قضیه، مثال، اثبات و نتیجه بیان شده است. می توان از محیط flushright برای مطالب راست چین و از محیط flushleft برای مطالب چپ چین استفاده نمود تا نمایش مطالب دقیقتر باشد.
    \end{flushright}
    \pause
    \begin{theorem}
        قضیه
    \end{theorem} 
    \pause
    \begin{example}
        مثال
    \end{example}
    \pause
    \begin{proof}
        اثبات
    \end{proof}
    \pause
    \begin{corollary}
        نتیجه
    \end{corollary}
\end{frame}
%%%%%%%%%%%%%%%%%%%%%%%%%%%%%%%%%%%%%%%%%%%%%%%%%%%%%%%%%%%%
\begin{frame}[c]
\begin{lemma}
روش اول تعریف محیط لم
\end{lemma}
\end{frame}
%%%%%%%%%%%%%%%%%%%%%%%%%%%%%%%%%%%%%%%%%%%%%%%%%%%%%%%%%%%%
\begin{frame}
\begin{block}{لم}
روش دوم تعریف محیط لم
\end{block}
\end{frame}
%%%%%%%%%%%%%%%%%%%%%%%%%%%%%%%%%%%%%%%%%%%%%%%%%%%%%%%%%%%%
\begin{frame}
\frametitle{Maths Blocks}
\begin{theorem}<1->[Pythagoras] 
$ a^2 + b^2 = c^2$
\end{theorem}
\begin{proof}<2->
$\omega +\phi = \epsilon $
\end{proof}
\begin{corollary}<3->
$ x + y = y + x  $
\end{corollary}
\end{frame}
%%%%%%%%%%%%%%%%%%%%%%%%%%%%%%%%%%%%%%%%%%%%%%%%%%%%%%%%%%%%%%
\subsection{بلوک ها}
%%%%%%%%%%%%%%%%%%%%%%%%%%%%%%%%%%%%%%%%%%%%%%%%%%%%%%%%%%%%
\begin{frame}[t]{انواع بلوک ها در بیمر}
با تایپ دستور
\lr{texdoc beamer}
در محیط کامند ویندوز یا ترمینال لینوکس می توانید به راهنمای بیمر دسترسی پیدا کنید.
همچنین میتوانید به آدرس
\lr{https://ctan.org/pkg/beamer?lang=en}
مراجعه بفرمایید.
\begin{block}{محیط بلوک معمولی}
متن بلوک معمولی
\end{block}
\begin{exampleblock}{محیط بلوک مثال}
متن بلوک مثال
\end{exampleblock}
\begin{alertblock}{محیط بلوک اعلان}
متن بلوک اعلان
\end{alertblock}
\end{frame}
%%%%%%%%%%%%%%%%%%%%%%%%%%%%%%%%%%%%%%%%%%%%%%%%%%%%%%%%%%%%
\subsection{محیط های شمارش و آیتم ها}
%%%%%%%%%%%%%%%%%%%%%%%%%%%%%%%%%%%%%%%%%%%%%%%%%%%%%%%%%%%%%%
\begin{frame}[t]{محیط شمارش- روش اول نمایش گزینه ها}
\begin{block}{شمارش}
    \begin{enumerate}
        \item گزینه اول \pause
        \item گزینه دوم \pause
        \item گزینه سوم \pause
        \item گزینه چهارم \pause
        \item گزینه پنجم \pause
    \end{enumerate}   
\end{block}	
\end{frame}
%%%%%%%%%%%%%%%%%%%%%%%%%%%%%%%%%%%%%%%%%%%%%%%%%%%%%%%%%%%%
\begin{frame}[t]{محیط شمارش- روش دوم نمایش گزینه ها}
\begin{block}{شمارش}
    \begin{enumerate}[<+->]
        \item گزینه اول 
        \item گزینه دوم 
        \item گزینه سوم 
        \item گزینه چهارم 
        \item گزینه پنجم 
    \end{enumerate}   
\end{block}	
\end{frame}
%%%%%%%%%%%%%%%%%%%%%%%%%%%%%%%%%%%%%%%%%%%%%%%%%%%%%%%%%%%%%
\begin{frame}[t]{محیط شمارش- بهترین روش نمایش گزینه ها}
این روش نمایش اسلاید بهترین روش برای نمایش گزینه ها به صورت تک تکی است.
از محیط flushright برای مطالب راست چین، مانند فارسی استفاده شود.
از محیط flushleft برای مطالب چپ چین، مانند انگلیسی استفاده شود.
\pause
\begin{block}{شمارش}
    \begin{enumerate}[<+->]
        \item \begin{flushright}
           گزینه اول  
        \end{flushright}
        \item \begin{flushright}
           گزینه دوم  
        \end{flushright}
        \item \begin{flushright}
           گزینه سوم  
        \end{flushright} 
        \item \begin{flushright}
           گزینه چهارم  
        \end{flushright}
        \item \begin{flushright}
           گزینه پنجم
        \end{flushright} 
    \end{enumerate}   
\end{block}	
\end{frame}
%%%%%%%%%%%%%%%%%%%%%%%%%%%%%%%%%%%%%%%%%%%%%%%%%%%%%%%%%%%%%
\begin{frame}[t]{محیط شمارش تو در تو}
\begin{block}{شمارش}
    \begin{enumerate}
        \item گزینه اول
        \item گزینه دوم
        \begin{enumerate}
            \item گزینه اول
            \item گزینه دوم
        \end{enumerate} 
    \end{enumerate}   
\end{block}	
\end{frame}
%%%%%%%%%%%%%%%%%%%%%%%%%%%%%%%%%%%%%%%%%%%%%%%%%%%%%%%%%%%%%
\begin{frame}[t]{محیط آیتم های تو در تو}
\begin{block}{آیتم ها}
    \begin{itemize}
        \item گزینه اول
        \item گزینه دوم
        \begin{itemize}
            \item گزینه اول
            \item گزینه دوم
        \end{itemize} 
    \end{itemize}   
\end{block}	
\end{frame}
%%%%%%%%%%%%%%%%%%%%%%%%%%%%%%%%%%%%%%%%%%%%%%%%%%%%%%%%%%%
\begin{frame}[t]{\vspace*{-8mm}\begin{flushleft}More Lists\end{flushleft}\vspace*{-8mm}}
\begin{latin}
\begin{enumerate}[(I)]
\item<1-> Point A
\item<2-> Point B
\begin{itemize}
\item<3-> part 1
\item<4-> part 2
\end{itemize}
\item<5-> Point C
\item<6-> Point D
\end{enumerate}
\end{latin}
\end{frame}
%%%%%%%%%%%%%%%%%%%%%%%%%%%%%%%%%%%%%%%%%%%%%%%%%%%%%%%%%%%%%
\begin{frame}{اسلاید دو ستونی}
    \begin{columns} 
     \column{.5\textwidth}
    ستون شماره ۱
     \column{.5\textwidth}
    ستون شماره ۲
    \end{columns}
\end{frame}
%%%%%%%%%%%%%%%%%%%%%%%%%%%%%%%%%%%%%%%%%%%%%%%%%%%%%%%%%%%%%%%
\setbeamercovered{invisible}
\begin{frame}
\frametitle{Tables}
\begin{table}
\begin{tabular}{l | c | c | c | c }
Competitor Name & Swim & Cycle & Run & Total \\
\hline \hline
John T & 13:04 & 24:15 & 18:34 & 55:53 \onslide<2-> \\ 
Norman P & 8:00 & 22:45 & 23:02 & 53:47 \onslide<3->\\
Alex K & 14:00 & 28:00 & n/a & n/a \onslide<4->\\
Sarah H & 9:22 & 21:10 & 24:03 & 54:35 
\end{tabular}
\caption{Triathlon results}
\end{table}
\end{frame}
%%%%%%%%%%%%%%%%%%%%%%%%%%%%%%%%%%%%%%%%%%%%%%%%%%%%%%%%%%%%%%%%
\section{سه مفهوم پایه: زبان، گرامر، اتوماتا}
\subsection{اتوماتا:}
\begin{frame}{اتوماتا:}
\begin{itemize}[<+->]
\item اتوماتا: جمع اتوماتون 
\vfill

\item اتوماتا: یک مدل انتزاعی از کامپیوتر
\vfill 

\item  اتوماتا دستگاه یا سیستمی است که برای ورودی‌های مشخص، کارکرد مشخصی انجام می‌دهد.
\end{itemize}

\pause
\vfill
\begin{block}{اتوماتا}
یک اتوماتا از بخش‌های زیر تشکیل شده است:

\begin{itemize}[<+->]
\item یک نوار ورودی که می‌تواند یک رشته را نگهداری کند با هد خواندن  (واحد ورودی)

\item یک نوار خروجی با هد نوشتن (واحد خروجی)

\item  یک حافظه (واحد حافظه)

\item  واحد کنترل (واحد محاسبه و منطق)
\end{itemize}
\end{block}
\end{frame}
%%%%%%%%%%%%%%%%%%%%%%%%%%%%%%%%%%%%%%%%%%%%%%%%%%%%%%%%%%%%%%%%%%%
\begin{frame}{اتوماتا:}
\begin{block}{اتوماتا}
یک اتوماتا از بخش‌های زیر تشکیل شده است:

\begin{itemize}
\item یک نوار ورودی که می‌تواند یک رشته را نگهداری کند با هد خواندن  (واحد ورودی)

\item یک نوار خروجی با هد نوشتن (واحد خروجی)

\item  یک حافظه (واحد حافظه)

\item  واحد کنترل (واحد محاسبه و منطق)
\end{itemize}
\end{block}
\centerline{\includegraphics[width=0.4\textwidth]{automata}}
\end{frame}
%%%%%%%%%%%%%%%%%%%%%%%%%%%%%%%%%%%%%%%%%%%%%%%%%%%%%%%%%%%%%%%%%%%%
{\usebackgroundtemplate{\includegraphics[width=\paperwidth,height=\paperheight]{back5.jpg}}

\begin{frame}[plain,noframenumbering]
\transglitter[duration=2]
%\vspace*{2.5cm}
\begin{center}
\Huge {\nas
\textcolor{red}{ با آرزوی }

\hspace*{2.5cm}\textcolor{green}{موفقیت}

\hspace*{5cm}\textcolor{red}{ و }

\hspace*{6cm}\textcolor{yellow}{ بهروزی}
}
\end{center}
\end{frame}
}
%%%%%%%%%%%%%%%%%%%%%%%%%%%%%%%%%%%%%%%%%%%%%%%%%%%%%%%%%%%%%%%%
\end{document}